\begin{proof}[Proof of Theorem \ref{thm:fair-trees}]
It is obvious that when $P$ and $Q$ are isomorphic, then $K_P=K_Q$.
Let us prove the converse.

We shall use induction on the size of $P$, with the case of size 1 being trivial.

Assuming now that the size of $P$ is at least $2$, 
we may decompose $P$ and $Q$ uniquely into non-empty
connected components:
$P=\sqcup_{i=1}^r P_i$
and
$Q=\sqcup_{i=1}^s Q_i$.
We have
\[
\prod_{i=1}^r K_{P_i}
=
\prod_{i=1}^s K_{Q_i}.
\]
Since \qsym\  is a unique factorization domain \cite{Haz01,LaPy08,MaRe95}, 
the irreducibility of $K_{P_i}$ and $K_{Q_i}$ from Proposition~\ref{prop:irred-C} 
implies that $r=s$ and that for every $i$, we have $K_{P_i} = K_{Q_i}$ (up to a suitable renumbering).

When $r \ge 2$, the size of each $P_i$ and $Q_i$ is smaller than the size of $P$
and thus $P_i$ and $Q_i$ are isomorphic for every $i$ by the induction hypothesis. 
Thus $P$ and $Q$ are also isomorphic. 

Suppose now that $r = 1$, i.e. $P$ and $Q$ are connected. Thus, 
since their size is greater than $1$, they may be written as
$P=[1]\wP P'$ or $P=[1]\sP P'$
or
$P=P'\wP [1]$ or $P=P' \sP [1]$,
and similarly for $Q$.
By Theorem~\ref{theo:KinF}, we can distinguish among these four possibilities by observing whether the first part of the compositions $\alpha$ in the $F$-support of $K_P$ (equivalently $K_Q$) are all equal to $1$ or all greater than $1$, and similarly for the last part of all such $\alpha$. Specifically, $P=[1]\wP P'$ if and only if $\alpha_1 > 1$ for all $\alpha$, implying $Q=[1]\wP Q'$.  Similarly, $P=[1]\sP P'$ if and only if $\alpha_1=1$ for all $\alpha$, implying $Q=[1]\sP Q'$.  And $P=P'\wP [1]$ (\resp $P=P' \sP [1]$) if and only if the last part of $\alpha$ is greater than 1 (\resp equals 1) for all $\alpha$, implying $Q=Q'\wP [1]$ (\resp $Q=Q' \sP [1]$).  
So for a given $P$ and $Q$ with $K_P = K_Q$\,, at least one of the following four properties holds:
\begin{itemize}
\item $P=[1]\wP P'$ and $Q=[1]\wP Q'$, or 
\item $P=[1]\sP P'$ and $Q=[1]\sP Q'$, or
\item $P=P'\wP [1]$ and $Q=Q'\wP [1]$, or 
\item $P=P'\sP [1]$ and $Q=Q'\sP [1]$.
\end{itemize}

In the first case, applying Proposition~\ref{prop:k_products} gives $F_1\wP K_{P'} = F_1\wP K_{Q'}$.  Let us expand this in the $F$-basis as $\sum c_\alpha F_\alpha$.  From~\eqref{equ:f_uparrow} and~\eqref{equ:uparrow}, we see that by just subtracting $1$ from the first part $\alpha_1$ in each term $F_\alpha$, we get $K_{P'} = K_{Q'}$, with $P'$ of size smaller than $P$.
By induction, $P'$ and $Q'$ are isomorphic, and hence so are $P$ and $Q$.

In the second case, we have  $F_1\sP K_{P'} = F_1\sP K_{Q'}$.
Let us expand this in the $F$-basis as $\sum c_\alpha F_\alpha$.
From~\eqref{equ:f_upuparrow} and~\eqref{equ:upuparrow}, we see that by just  removing the first part $\alpha_1=1$ 
in each term $F_\alpha$, we get $K_{P'} = K_{Q'}$, with $P'$ of size smaller than $P$, giving the same conclusion.

The last two cases are resolved similarly by focusing on the last entry of $\alpha$.
\end{proof}
